\begin{obeactivity}
\textbf{Aktivitas 13.1: Simulasi Merkle-Hellman}
1. Tentukan kunci privat (barisan \textit{superincreasing}): $W_{priv} = \{2, 3,
   6, 13, 27, 52\}$.
2. Gunakan $m = 105$ dan $n = 31$. Periksa lebih dahulu bahwa $m > \sum_i w_i$
   dan bahwa $\text{PBB}(n, m) = 1$. Mengapa kedua syarat itu wajib?
3. Hitung kunci publik $W_{pub}$ menurut Persamaan~\eqref{eq:mh-kunci}.
4. Bandingkan $W_{priv}$ dengan $W_{pub}$. Tunjukkan dengan memeriksa setiap
   elemennya bahwa $W_{pub}$ bukan lagi barisan superincreasing.
5. Hitung $n^{-1} \bmod 105$.
\end{obeactivity}

Dua perkakas pada Listing~\ref{lst:knapsack-kunci} menjalankan langkah 3 dan 5
di atas: memotong rangkaian bit menjadi blok sepanjang $n$ bit, dan menghitung
barisan publik dari barisan privat.

\lstinputlisting[
    language=Python,
    caption={Pemotongan blok dan pembangkitan barisan publik Merkle-Hellman},
    label={lst:knapsack-kunci},
    firstline=113,
    lastline=123
]{\codefile{python/knapsack_uji.py}}

\begin{obeactivity}
\textbf{Aktivitas 13.2: Dekripsi \textit{Greedy} pada Barisan Superincreasing}
1. Periksa bahwa barisan $\{2, 3, 6, 13, 27, 52\}$ memenuhi
   Persamaan~\eqref{eq:superincreasing} untuk setiap elemennya.
2. Jalankan fungsi pada Listing~\ref{lst:knapsack-greedy} untuk memecahkan
   $M = 70$.
3. Tuliskan langkah demi langkahnya seperti pada
   Tabel~\ref{tab:greedy-70}, dan bandingkan hasilnya dengan
   Contoh~\ref{ex:greedy-superincreasing}.
4. Ubah $M$ menjadi $71$ dan jalankan kembali fungsinya. Mengapa fungsi itu
   mengembalikan \emph{tidak ada solusi}, dan apa artinya bagi keamanan pesan?
5. Ulangi untuk $M = 103$ (jumlah seluruh elemen). Bit apa yang dihasilkan?
\end{obeactivity}

\lstinputlisting[
    language=Python,
    caption={Pemeriksaan barisan superincreasing dan dekripsi \textit{greedy}},
    label={lst:knapsack-greedy},
    firstline=64,
    lastline=82
]{\codefile{python/knapsack_uji.py}}

\begin{obeactivity}
\textbf{Aktivitas 13.3: Enkripsi dan Dekripsi Tiga Blok}
1. Jalankan program pada Listing~\ref{lst:knapsack-enkripsi}, yang mengenkripsi
   \code{011000110101101110} dengan $W_{pub} = \{62, 93, 81, 88, 102, 37\}$.
2. Periksa bahwa kriptogramnya $(174, 280, 333)$.
3. Perhatikan bahwa $280$ dan $333$ melampaui $m = 105$. Jelaskan mengapa hal
   itu tidak merusak dekripsi.
4. Hitung $C' = (C \cdot n^{-1}) \bmod 105$ untuk ketiga kriptogram, lalu
   pecahkan hasilnya dengan algoritma \textit{greedy} pada $W_{priv}$.
5. Bandingkan hasilnya dengan Tabel~\ref{tab:mh-105-dekripsi}. Apakah ketiga
   blok pulih tepat seperti semula?
\end{obeactivity}

\lstinputlisting[
    language=Python,
    caption={Enkripsi dan dekripsi tiga blok dengan $m = 105$ dan $n = 31$},
    label={lst:knapsack-enkripsi},
    firstline=186,
    lastline=207
]{\codefile{python/knapsack_uji.py}}

\begin{obeactivity}
\textbf{Aktivitas 13.4: Mengukur Kerapatan dan Memeriksa Parameter}
1. Hitung kerapatan barisan $\{62, 93, 81, 88, 102, 37\}$ dengan
   Persamaan~\eqref{eq:densitas-knapsack}.
2. Hitung kerapatan barisan Merkle--Hellman yang dianjurkan sumber: 100 elemen
   dengan elemen terbesar sekitar 299 bit.
3. Bandingkan kedua nilai itu dengan ambang $d < 0{,}645$
   (Lagarias--Odlyzko) dan $d < 0{,}9408$ (Coster dan kawan-kawan).
   Barisan manakah yang justru dapat diserang?
4. Periksa anjuran modulus pada sumber: bisakah $m$ sepanjang 100--200 bit
   melampaui jumlah 250 elemen yang masing-masing bernilai 200 bit ke atas?
5. Bangkitkan barisan superincreasing sebanyak 100 elemen mulai dari sekitar
   $2^{200}$ dan hitung panjang bit jumlahnya. Berapa panjang $m$ yang
   sesungguhnya diperlukan?
\end{obeactivity}

\lstinputlisting[
    language=Python,
    caption={Kerapatan barisan, biaya penelusuran, dan pemeriksaan parameter},
    label={lst:knapsack-kerapatan},
    firstline=242,
    lastline=289
]{\codefile{python/knapsack_uji.py}}

\begin{obeactivity}
\textbf{Aktivitas 13.5: Membalik Kunci Publik Menjadi Kunci Privat}
1. Jalankan program pada Listing~\ref{lst:knapsack-balik}.
2. Perhatikan bahwa mengalikan $W_{pub}$ dengan $n^{-1}$ modulo $m$
   mengembalikan $W_{priv}$ secara tepat.
3. Jelaskan mengapa langkah itu sama persis dengan langkah kedua pada prosedur
   dekripsi (Subbagian~\ref{subsec:bukti-knapsack}).
4. Sekarang pikirkan posisi penyerang. Bila ia mengetahui barisan publik tetapi
   bukan $(n, m)$, persoalan apa yang harus ia pecahkan? Berapa banyak pasangan
   $(n, m)$ yang mungkin, dan mengapa menelusuri seluruhnya tidak praktis?
5. Uraikan secara ringkas mengapa Shamir dapat menemukan $(n, m)$ itu dalam
   waktu polinomial, sedangkan penelusuran menyeluruh tidak mungkin
   (Subbagian~\ref{sec:kriptanalisis-knapsack}).
\end{obeactivity}

\lstinputlisting[
    language=Python,
    caption={Pemulihan kunci privat dari kunci publik},
    label={lst:knapsack-balik},
    firstline=292,
    lastline=303
]{\codefile{python/knapsack_uji.py}}

Seluruh sepuluh pengujian pada program ini dapat dijalankan sekaligus dengan
\texttt{python knapsack\_uji.py}. Bila setiap angka pada bab ini cocok, program
akan mencetak \texttt{LULUS} pada setiap blok pengujian.
